%----------------------------------------------------------------------------------------
%   Түлхүүр үгийн тайлбар
%   Шинэ үг нэмэхдээ жагсаалтад \item нэмнэ. Цагаан толгойн дарааллыг хадгална уу.
%----------------------------------------------------------------------------------------
\newpage
\begingroup
\renewcommand{\headrulewidth}{0pt}
\markboth{}{}
\begin{center}
\textbf{\large ТҮЛХҮҮР ҮГИЙН ТАЙЛБАР}
\end{center}

\singlespace

\noindent\textbf{Товчилсон үг}

\begin{description}
\item[2D, 3D] Two-, Three-Dimensional --- Хоёр хэмжээст, гурван хэмжээст
\item[API] Application Programming Interface --- Програмууд хоорондоо хүсэлт илгээж харилцах интерфейс
\item[COCO] Common Objects in Context --- Хүний гол цэгийн тэмдэглэгээтэй, өргөн хэрэглэгддэг зургийн өгөгдлийн сан
\item[CPU] Central Processing Unit --- Компьютерын төв процессор
\item[CSV] Comma-Separated Values --- Таслалаар тусгаарласан хүснэгтэн өгөгдлийн файлын формат
\item[CUDA] Compute Unified Device Architecture --- NVIDIA-ийн GPU дээр ерөнхий тооцоолол хийх платформ
\item[CV] Coefficient of Variation --- Хэлбэлзлийн коэффициент
\item[DTW] Dynamic Time Warping --- Динамик цагийн хазайлт; хурд нь ялгаатай хоёр цувааг харьцуулах арга
\item[GPU] Graphics Processing Unit --- График процессор (видео карт)
\item[HPE] Human Pose Estimation --- Хүний байрлал тодорхойлох
\item[HRNet] High-Resolution Network --- Өндөр нягтаршилтай дүрслэлд суурилсан HPE загвар
\item[JSON] JavaScript Object Notation --- Бүтэцтэй өгөгдөл хадгалах текст формат
\item[LSP] Leeds Sports Pose --- Спортын байрлалын өгөгдлийн сан
\item[MAE] Mean Absolute Error --- Дундаж абсолют алдаа
\item[MP4] MPEG-4 Part 14 --- Видео файлын формат
\item[MPII] MPII Human Pose --- Хүний байрлалын өгөгдлийн сан
\item[NMS] Non-Maximum Suppression --- Давхацсан илрүүлэлтүүдээс хамгийн өндөр итгэлцэлтэйг үлдээх алхам
\item[ONNX] Open Neural Network Exchange --- Нейрон сүлжээний загварыг өөр орчинд шилжүүлэх нээлттэй формат
\item[PAF] Part Affinity Fields --- Хэсгийн холбоосын талбар; OpenPose-д гол цэгүүдийг хүн тус бүрээр холбоход ашиглана
\item[PCK] Percentage of Correct Keypoints --- Зөв илэрсэн гол цэгийн хувь
\item[RTMPose] Real-Time Multi-person Pose Estimation --- Бодит хугацааны олон хүний байрлал тодорхойлох загвар
\item[SHA-256] Secure Hash Algorithm (256-bit) --- Файл ижил эсэхийг шалгах хэш функц
\item[SimCC] Simple Coordinate Classification --- RTMPose-ийн гол цэгийн x, y координатыг тэнхлэг тус бүрээр ангилж тодорхойлох гаралт
\item[ViT] Vision Transformer --- Зургийн трансформер; ViTPose загварын суурь архитектур
\item[WASM] WebAssembly --- Вэб хөтөч дотор CPU дээр хурдан ажиллах кодын формат
\item[WebGL] Web Graphics Library --- Вэб хөтөчөөс GPU ашиглан график зурах технологи
\item[WebGPU]Вэб хөтөчөөс GPU-ийн тооцооллыг шууд ашиглах шинэ технологи
\item[YOLO] You Only Look Once --- Объект илрүүлэлтийг нэг алхамд гүйцэтгэдэг загваруудын гэр бүл
\end{description}

\noindent\textbf{Нэр томьёо}

\begin{description}
\item[Араг яс (skeleton)] Гол цэгүүдийг хооронд нь шугамаар холбосон хүний биеийн энгийн дүрслэл
\item[Гол цэг (keypoint)] Хамар, мөр, тохой, өвдөг зэрэг хүний биеийн анатомийн тодорхой цэг
\item[Гулсах цонх (sliding window)] Хугацааны цувааны тодорхой урттай хэсгийг шилжүүлэн тооцоолол хийх арга
\item[Дээрээс доош (top-down)] Эхлээд хүн бүрийг хүрээгээр илрүүлж, дараа нь хүрээ тус бүрт гол цэгийг тодорхойлдог HPE арга
\item[Доороос дээш (bottom-up)] Эхлээд бүх гол цэгийг илрүүлж, дараа нь хүн тус бүрээр бүлэглэдэг HPE арга
\item[Жинхэнэ утга (ground truth)] Гараар тэмдэглэсэн, зөв гэж тооцогдох жишиг өгөгдөл
\item[Консенсус (consensus)] Бусад загваруудын гаргасан байрлалын медиан; жинхэнэ утга байхгүй үед жишиг болгон ашиглана
\item[Корреляц (correlation)] Хоёр цуваа хэр зэрэг хамт өөрчлөгдөж байгааг $-1$-ээс 1 хүртэлх утгаар илэрхийлэх хэмжүүр
\item[Сервер (server)] Хэрэглэгчийн хүсэлтээр тооцоолол гүйцэтгэж үр дүнг буцаадаг алсын компьютер; энэ ажилд видеоны гол цэгийг тодорхойлно
\item[Синхрончлол (synchronization)] Хоёр болон түүнээс дээш хүний хөдөлгөөн хэлбэр, хэмнэл, цаг хугацааны хувьд хэр нийцэж буй байдал
\item[Тооцооллын backend (execution provider)] Загварын тооцооллыг гүйцэтгэх технологи (WebGPU, WASM, CUDA гэх мэт)
\item[Үений өнцөг (joint angle)] Гурван гол цэгээр тодорхойлогдох үений нугаралтын өнцөг, жишээлбэл тохойн өнцөг
\item[Хөндлөн корреляц (cross-correlation)] Нэг цувааг цаг хугацааны дагуу шилжүүлэн нөгөөтэй нь харьцуулж, цагийн хоцрогдлыг олох арга
\item[Хяналттай өгөгдөл (synthetic data)] Хоцрогдол, зай зэрэг ялгааг судлаач өөрөө тогтоож үүсгэсэн туршилтын өгөгдөл
\item[Цагийн хоцрогдол (lag)] Нэг хүний хөдөлгөөн нөгөөгөөсөө хэдэн кадр эсвэл секундээр хоцорч байгааг илэрхийлэх утга
\item[Чичиргээ (jitter)] Хүн хөдөлгөөнгүй байсан ч дараалсан кадруудад гол цэгийн байрлал санамсаргүй үсрэх үзэгдэл
\item[Шуугиан тэгшитгэх (smoothing)] Хөрш кадруудын утгыг ашиглан чичиргээг багасгаж, хөдөлгөөний муруйг гөлгөр болгох алхам
\end{description}

\doublespace
\newpage
\endgroup
